% Layout-pattern reference for latex-presentation decks (spec:
% spec-latex-presentation-generation.md). Extracted from the researcher's
% reference template's example body frames.
%
% These are STYLE EXAMPLES of layout structure -- columns, equations,
% block+figure -- not content to reproduce verbatim in a generated deck.
% Draw the *pattern* (how a frame is laid out) and drop in the researcher's
% actual content; never copy the placeholder text/labels below into a real
% presentation. Not a standalone compilable file -- these frames are meant
% to be lifted into a presentations/<slug>/main.tex that already has
% default-preamble.tex's preamble (or an alternate one, on a theming
% override) in place.

% --- Pattern: two-column text (itemize) + figure -----------------------
% Use for: introducing a topic with a short bulleted summary alongside a
% supporting image.
\begin{frame}
\frametitle{Title slide}
\begin{minipage}{\textwidth}
\begin{columns}
\begin{column}{0.5\textwidth}
Normal text
\begin{itemize}
    \item item text
    \item item text 2
    \item item text 3
\end{itemize}
\end{column}
\begin{column}{0.5\textwidth}
\begin{figure}
\includegraphics[width=6.5cm]{images/graph_social_media.png}
\caption{Title image}
\end{figure}
\end{column}
\end{columns}
\end{minipage}
\end{frame}

% --- Pattern: equations + itemize (math-heavy slide) --------------------
% Use for: defining notation/formalism before results depend on it.
\begin{frame}
\frametitle{Some equations}
$G(V, E)$
\begin{itemize}
    \item $V$ super pixels
    \item $E$ set of edges
    \item $d$  number of connections
\end{itemize}
Matrix of embeddings $H$:
\[ \mathcal H^{(l)} = [h_1^{(l)} ... h_{\left | V \right |}^{(l)}]^{\text{T}}\]
\begin{itemize}
    \item $|V|$ number of nodes
    \item $l$ model´s layer
\end{itemize}
\end{frame}

% --- Pattern: block + figure side-by-side -------------------------------
% Use for: pairing a highlighted explanatory block (definition, takeaway,
% caveat) with a supporting figure in the other column.
\begin{frame}
\frametitle{Slide with different columns}

\begin{columns}
\begin{column}{0.5\textwidth}

\begin{block}{A block}
Text inside my block. a block is kinda a slide inside a slide, in a rectangular shape.
\end{block}
\end{column}
\begin{column}{0.5\textwidth}
\begin{figure}
\includegraphics{images/GNN_layers.PNG}
\caption{an image on the right column}
\end{figure}
\end{column}
\end{columns}
\end{frame}
